\documentclass[aspectratio=169,10pt]{beamer}

% --- Paquetes Básicos ---
\usepackage[utf8]{inputenc}
\usepackage[spanish]{babel}
\usepackage{graphicx}
\usepackage{booktabs}
\usepackage{amsmath}
\usepackage{amsfonts}
\usepackage{tikz}
\usetikzlibrary{shapes.geometric, arrows, positioning, fit, backgrounds}

% --- Tema y Colores Institucionales UTEC Posgrado ---
\usetheme{Madrid}
\usecolortheme{whale}

% Colores de la plantilla UTEC Posgrado (Gris oscuro y Azul cian)
\definecolor{UtecCyan}{RGB}{0, 191, 255}       % Azul cian brillante
\definecolor{UtecDarkGray}{RGB}{50, 50, 50}    % Gris Oscuro institucional
\definecolor{LightGray}{RGB}{245, 245, 245}    % Fondo suave
\definecolor{White}{RGB}{255, 255, 255}

% Aplicar colores de marca a Beamer
\setbeamercolor{palette primary}{bg=UtecDarkGray, fg=white}
\setbeamercolor{palette secondary}{bg=UtecDarkGray!85, fg=white}
\setbeamercolor{palette tertiary}{bg=UtecCyan, fg=UtecDarkGray}
\setbeamercolor{titlelike}{parent=palette primary, fg=UtecDarkGray}
\setbeamercolor{structure}{fg=UtecCyan}
\setbeamercolor{background canvas}{bg=white}
\setbeamercolor{normal text}{fg=UtecDarkGray}

% Estilo de bloques beamer
\setbeamercolor{block title}{bg=UtecDarkGray, fg=white}
\setbeamercolor{block body}{bg=LightGray, fg=UtecDarkGray}

% Quitar barra de navegación de Beamer para mayor limpieza visual
\setbeamertemplate{navigation symbols}{}

% Pie de página minimalista
\setbeamertemplate{footline}[frame number]

% --- Metadatos de la Presentación ---
\title[Ingeniería de Datos - Proyecto Final]{INGENIERÍA DE DATOS}
\subtitle{Proyecto Final: Medallion Architecture \& Data Mesh}
\author[Grupo G2]{Integrantes: \\ \textbf{Grupo G2 (Sección 1 - g102)}}
\institute[UTEC]{Universidad de Ingeniería y Tecnología - UTEC Posgrado}
\date{18 de Julio de 2026}

\begin{document}

% --- Slide 1: Portada (Plantilla UTEC) ---
{
\setbeamercolor{background canvas}{bg=UtecDarkGray}
\begin{frame}[plain]
    \centering
    \vspace{1.5cm}
    {\large \color{UtecCyan} \textbf{CURSO}} \\[0.2cm]
    {\Huge \color{white} \textbf{INGENIERÍA DE DATOS}} \\[0.6cm]
    \color{white!80}
    {\large \textbf{Proyecto Final - Grupo G2}} \\[1.2cm]
    \footnotesize
    \color{white!70}
    Integrantes: herles.pinedo@utec.edu.pe | david.jimenez@utec.edu.pe | jose.nomberto@utec.edu.pe \\
    kevin.luyo@utec.edu.pe | yessenia.chirinos@utec.edu.pe \\
    \vspace{0.4cm}
    UTEC Posgrado | Julio 2026
\end{frame}
}

% --- Slide 2: Nombre de la Empresa y Rubro (Plantilla UTEC) ---
\begin{frame}{Contexto Organizacional}
    \begin{columns}
        \begin{column}{0.5\textwidth}
            {\Huge \color{UtecDarkGray} \textbf{Compañía Minera Los Andes S.A.}} \\
            \vspace{0.6cm}
            {\Large \color{UtecCyan} \textbf{RUBRO: Minería polimetálica}} \\
            \vspace{0.6cm}
            \small
            \textbf{Contexto:} Operación a tajo abierto con alta dependencia en la disponibilidad y consumo de repuestos críticos bajo esquemas de inventario en consignación en mina.
        \end{column}
        \begin{column}{0.45\textwidth}
            \centering
            \begin{tikzpicture}[scale=0.8, every node/.style={scale=0.8}]
                \draw[fill=LightGray, draw=UtecDarkGray, thick] (0,0) rectangle (4,3);
                \node at (2,2.5) [text=UtecDarkGray, font=\bfseries\small] {Operación Minera};
                \draw[fill=UtecCyan!30, draw=UtecDarkGray] (0.5,0.5) rectangle (1.5,1.5);
                \draw[fill=UtecDarkGray!20, draw=UtecDarkGray] (2.5,0.5) rectangle (3.5,1.5);
                \node at (1,1) [text=UtecDarkGray, scale=0.6, align=center] {Tajo\\Abierto};
                \node at (3,1) [text=UtecDarkGray, scale=0.6, align=center] {Planta\\Concentradora};
                \draw[->, thick, draw=UtecDarkGray] (1.5,1) -- (2.5,1);
            \end{tikzpicture}
        \end{column}
    \end{columns}
\end{frame}

% --- Slide 3: Ejemplos - Datasets por rubro (Plantilla UTEC) ---
\begin{frame}{Ejemplos - Datasets por rubro}
    \centering
    \small
    \begin{table}[h]
        \begin{tabular}{lll}
            \toprule
            \textbf{Rubro} & \textbf{Dataset} & \textbf{Tipos de datos} \\
            \midrule
            Retail & sales.csv, products.csv & Transaccional \\
            Salud & appointments.csv, patients.csv & Registros clínicos \\
            Finanzas & transactions.csv, customers.csv & Movimientos financieros \\
            \rowcolor{UtecCyan!20} \textbf{Logística (Nuestro Caso)} & \textbf{MB51\_*.csv, maestro\_materiales.csv} & \textbf{Movimientos e Inventario SAP} \\
            \bottomrule
        \end{tabular}
    \end{table}
    \vspace{0.4cm}
    \footnotesize
    \color{UtecDarkGray!70}
    Nota: Para nuestro proyecto final, procesamos la data transaccional real de movimientos históricos de almacén de SAP del rubro logístico.
\end{frame}

% --- Slide 4: WHY (Plantilla UTEC) ---
\begin{frame}{WHY -- Por qué estamos haciendo esto?}
    \textbf{Objetivo: Justificar la existencia del Data Product desde un problema real del negocio.}
    \vspace{0.4cm}
    \begin{itemize}
        \item \textbf{¿Qué necesidad, dolor o desafío existe en el dominio del negocio?}
          Evitar el riesgo de desabastecimiento en mina y la parada de planta concentradora por falta de repuestos críticos en consignación (ej. medios de molienda y lubricantes).
        \item \textbf{¿Qué impacto tiene no resolver este problema?}
          Pérdidas millonarias por horas de parada de planta y exceso de capital de trabajo inmovilizado.
        \item \textbf{¿Por qué un Data Product es la mejor solución?}
          Centraliza la lógica en un producto gobernado, con SLAs de latencia y contratos de calidad de datos.
        \item \textbf{¿Qué valor generará?}
          Reducción del 95\% del tiempo de parada preventivo y automatización de la conciliación mensual de pagos.
    \end{itemize}
\end{frame}

% --- Slide 5: WHAT (Plantilla UTEC) ---
\begin{frame}{WHAT -- Qué estamos construyendo?}
    \textbf{Objetivo: Describir el Data Product como unidad de valor.}
    \vspace{0.4cm}
    \begin{itemize}
        \item \textbf{¿Cuál es el Data Product? ¿Qué representa y para quién?}
          Portal de alertas y costos de inventario de consignación, consumido por Planeamiento de Mantenimiento y Finanzas.
        \item \textbf{¿Qué dominio lo produce? ¿Quién lo consume?}
          Producido por el dominio de Logística y consumido por Finanzas y Abastecimiento.
        \item \textbf{¿Qué datos incluye y por qué son relevantes?}
          Alertas de stock de seguridad, montos de reposición en USD y transacciones pre-aprobadas listas para facturar.
        \item \textbf{¿Qué nivel de disponibilidad y calidad ofrece (SLOs)?}
          Disponibilidad diaria a primera hora, latencia de consulta $<2.5$ segundos y 100\% de integridad en IDs clave.
        \item \textbf{¿Cómo se ve el resultado final?}
          Vistas gobernadas en la capa Gold y un Dashboard interactivo con filtros por proveedor.
    \end{itemize}
\end{frame}

% --- Slide 6: HOW (Plantilla UTEC) ---
\begin{frame}{HOW -- Cómo lo estamos implementando?}
    \textbf{Objetivo: Mostrar cómo se construye y gobierna el Data Product.}
    \vspace{0.4cm}
    \begin{itemize}
        \item \textbf{¿Cómo se estructuró el pipeline (bronze/silver/gold)?}
          Bronze (ingesta inmutable y limpieza de nombres) ➡️ Silver (calidad, enmascaramiento UDF y cuarentena) ➡️ Gold (capa semántica de negocio).
        \item \textbf{¿Se usaron herramientas como Unity Catalog, contratos YAML, validaciones, workflows?}
          Sí, control de acceso y enmascaramiento dinámico PII en Unity Catalog, validaciones de calidad con contratos YAML y automatización diaria mediante Databricks Workflows.
        \item \textbf{¿Qué reglas de calidad o clasificación se implementaron?}
          Desvío automático de registros con cantidad nula, fechas incorrectas o materiales inválidos a Cuarentena.
        \item \textbf{¿Qué decisiones técnicas y de gobernanza se tomaron?}
          Uso de `try_cast` para evitar caídas de tareas por datos corruptos de SAP y vistas lógicas estándar en Gold.
    \end{itemize}
\end{frame}

% --- Slide 7: WHEN (Plantilla UTEC) ---
\begin{frame}{WHEN -- Cuál es el estado y próximos pasos?}
    \textbf{Objetivo: Mostrar el progreso actual y cómo se completará el proyecto.}
    \vspace{0.4cm}
    \begin{itemize}
        \item \textbf{¿Qué ya está hecho (Mínimo viable)?}
          Pipeline funcional en Databricks de Bronze a Gold y Dashboard interactivo en Databricks Lakeview en producción.
        \item \textbf{¿Qué falta implementar o validar?}
          Aprobación final de los catálogos en Unity Catalog por el área central de TI y conectar alertas a Teams.
        \item \textbf{¿Qué desafíos has identificado y cómo lo abordarás?}
          Estandarización de las exportaciones de SAP. Lo abordamos con reglas robustas de enmascaramiento y try-casts.
        \item \textbf{¿Qué plan tienes hasta la entrega final?}
          Validación del plan del Workflow con el CDO y carga de contratos definitivos en el repositorio GitHub antes del 18 de Julio.
    \end{itemize}
\end{frame}

% --- Slide 8: Recomendaciones Finales (Plantilla UTEC) ---
\begin{frame}{Recomendaciones Finales}
    \begin{itemize}
        \setlength{\itemsep}{0.6em}
        \item \textbf{Enfoque en Valor:} Priorizar siempre la claridad y el valor operativo al negocio por encima de la complejidad técnica de la nube.
        \item \textbf{Conceptos de Data Mesh:} Adoptar un lenguaje común de dominios, contratos de datos, productos y consumo federado.
        \item \textbf{Trazabilidad y Calidad:} Mantener la auditoría de SLAs y el control del ciclo de vida del dato mediante contratos YAML.
        \item \textbf{Hacerlo Simple y Gobernable:} Apoyar el consumo de los Data Products con visualizaciones amigables para el usuario final.
    \end{itemize}
    \vspace{0.3cm}
    \centering
    \small \textbf{“No se trata de procesar datos. Se trata de crear un producto de datos útil, gobernado y alineado al negocio.”}
\end{frame}

% --- Slide 9: Cierre (Plantilla UTEC) ---
{
\setbeamercolor{background canvas}{bg=UtecDarkGray}
\begin{frame}[plain]
    \centering
    \vspace{2cm}
    {\Huge \color{UtecCyan} \textbf{¡Gracias!}} \\[0.8cm]
    \color{white!80}
    {\large \textbf{Preguntas o comentarios}} \\[1cm]
    \footnotesize
    \color{white!60}
    Curso: Ingeniería de Datos | UTEC Posgrado
\end{frame}
}

\end{document}
